\documentclass[11pt]{article}
% Preambulo comun para los informes de la ronda 3.
% Se incluye con % Preambulo comun para los informes de la ronda 3.
% Se incluye con % Preambulo comun para los informes de la ronda 3.
% Se incluye con \input{../comun/preambulo} desde cada subcarpeta.
\usepackage[utf8]{inputenc}
\usepackage[T1]{fontenc}
\usepackage[spanish,es-nodecimaldot]{babel}
\usepackage{lmodern}
\usepackage[a4paper,margin=2.4cm]{geometry}
\usepackage{booktabs}
\usepackage{longtable}
\usepackage{array}
\usepackage{xcolor}
\usepackage{hyperref}
\usepackage{enumitem}
\usepackage{fancyhdr}
\usepackage{titlesec}
\usepackage{graphicx}
\usepackage{amssymb}
\usepackage{microtype}

\definecolor{sosTeal}{HTML}{0F766E}
\definecolor{sosGray}{HTML}{475569}

\hypersetup{
  colorlinks=true,
  linkcolor=sosTeal,
  urlcolor=sosTeal,
  citecolor=sosTeal,
  pdftitle={SOS Clinic - Ronda 3 de pruebas},
  pdfauthor={Equipo SOS Clinic}
}

\pagestyle{fancy}
\fancyhf{}
\lhead{\textcolor{sosGray}{SOS Clinic --- Ronda 3}}
\rhead{\textcolor{sosGray}{\leftmark}}
\cfoot{\thepage}
\setlength{\headheight}{15pt}
\renewcommand{\headrulewidth}{0.4pt}

\titleformat{\section}{\Large\bfseries\color{sosTeal}}{\thesection}{0.8em}{}
\titleformat{\subsection}{\large\bfseries\color{sosGray}}{\thesubsection}{0.6em}{}

\setlist[itemize]{leftmargin=1.4em,itemsep=2pt,topsep=4pt}
\setlist[enumerate]{leftmargin=1.6em,itemsep=2pt,topsep=4pt}

\newcommand{\pantalla}[1]{\texttt{\detokenize{#1}}}
\newcommand{\ok}{\textcolor{sosTeal}{\textbf{PASS}}}
\newcommand{\fallo}{\textcolor{red!70!black}{\textbf{FAIL}}}
\newcommand{\aviso}{\textcolor{orange!80!black}{\textbf{WARN}}}
\newcommand{\omitido}{\textcolor{sosGray}{\textbf{SKIP}}}
 desde cada subcarpeta.
\usepackage[utf8]{inputenc}
\usepackage[T1]{fontenc}
\usepackage[spanish,es-nodecimaldot]{babel}
\usepackage{lmodern}
\usepackage[a4paper,margin=2.4cm]{geometry}
\usepackage{booktabs}
\usepackage{longtable}
\usepackage{array}
\usepackage{xcolor}
\usepackage{hyperref}
\usepackage{enumitem}
\usepackage{fancyhdr}
\usepackage{titlesec}
\usepackage{graphicx}
\usepackage{amssymb}
\usepackage{microtype}

\definecolor{sosTeal}{HTML}{0F766E}
\definecolor{sosGray}{HTML}{475569}

\hypersetup{
  colorlinks=true,
  linkcolor=sosTeal,
  urlcolor=sosTeal,
  citecolor=sosTeal,
  pdftitle={SOS Clinic - Ronda 3 de pruebas},
  pdfauthor={Equipo SOS Clinic}
}

\pagestyle{fancy}
\fancyhf{}
\lhead{\textcolor{sosGray}{SOS Clinic --- Ronda 3}}
\rhead{\textcolor{sosGray}{\leftmark}}
\cfoot{\thepage}
\setlength{\headheight}{15pt}
\renewcommand{\headrulewidth}{0.4pt}

\titleformat{\section}{\Large\bfseries\color{sosTeal}}{\thesection}{0.8em}{}
\titleformat{\subsection}{\large\bfseries\color{sosGray}}{\thesubsection}{0.6em}{}

\setlist[itemize]{leftmargin=1.4em,itemsep=2pt,topsep=4pt}
\setlist[enumerate]{leftmargin=1.6em,itemsep=2pt,topsep=4pt}

\newcommand{\pantalla}[1]{\texttt{\detokenize{#1}}}
\newcommand{\ok}{\textcolor{sosTeal}{\textbf{PASS}}}
\newcommand{\fallo}{\textcolor{red!70!black}{\textbf{FAIL}}}
\newcommand{\aviso}{\textcolor{orange!80!black}{\textbf{WARN}}}
\newcommand{\omitido}{\textcolor{sosGray}{\textbf{SKIP}}}
 desde cada subcarpeta.
\usepackage[utf8]{inputenc}
\usepackage[T1]{fontenc}
\usepackage[spanish,es-nodecimaldot]{babel}
\usepackage{lmodern}
\usepackage[a4paper,margin=2.4cm]{geometry}
\usepackage{booktabs}
\usepackage{longtable}
\usepackage{array}
\usepackage{xcolor}
\usepackage{hyperref}
\usepackage{enumitem}
\usepackage{fancyhdr}
\usepackage{titlesec}
\usepackage{graphicx}
\usepackage{amssymb}
\usepackage{microtype}

\definecolor{sosTeal}{HTML}{0F766E}
\definecolor{sosGray}{HTML}{475569}

\hypersetup{
  colorlinks=true,
  linkcolor=sosTeal,
  urlcolor=sosTeal,
  citecolor=sosTeal,
  pdftitle={SOS Clinic - Ronda 3 de pruebas},
  pdfauthor={Equipo SOS Clinic}
}

\pagestyle{fancy}
\fancyhf{}
\lhead{\textcolor{sosGray}{SOS Clinic --- Ronda 3}}
\rhead{\textcolor{sosGray}{\leftmark}}
\cfoot{\thepage}
\setlength{\headheight}{15pt}
\renewcommand{\headrulewidth}{0.4pt}

\titleformat{\section}{\Large\bfseries\color{sosTeal}}{\thesection}{0.8em}{}
\titleformat{\subsection}{\large\bfseries\color{sosGray}}{\thesubsection}{0.6em}{}

\setlist[itemize]{leftmargin=1.4em,itemsep=2pt,topsep=4pt}
\setlist[enumerate]{leftmargin=1.6em,itemsep=2pt,topsep=4pt}

\newcommand{\pantalla}[1]{\texttt{\detokenize{#1}}}
\newcommand{\ok}{\textcolor{sosTeal}{\textbf{PASS}}}
\newcommand{\fallo}{\textcolor{red!70!black}{\textbf{FAIL}}}
\newcommand{\aviso}{\textcolor{orange!80!black}{\textbf{WARN}}}
\newcommand{\omitido}{\textcolor{sosGray}{\textbf{SKIP}}}


\title{\textbf{Pruebas Maven por módulo}\\[0.3em]
\large Metodología de ejecución --- Ronda 3}
\author{Equipo SOS Clinic}
\date{22 de octubre de 2026}

\begin{document}
\maketitle

\section{Objetivo}

Ejecutar la batería de pruebas existente en los siete proyectos Maven del
repositorio y dejar una constancia reproducible del resultado por módulo, sin
alterar el código ni la configuración de las pruebas.

\section{Alcance}

\begin{table}[h]
\centering
\small
\begin{tabular}{@{}ll@{}}
\toprule
Módulo & Tipo de proyecto \\
\midrule
\texttt{eureka-server} & Servidor de descubrimiento \\
\texttt{ms-gateway} & Gateway HTTP \\
\texttt{ms-seguridad} & Autenticación, JWT, 2FA, perfiles y turnos \\
\texttt{ms-pacientes} & Pacientes, signos vitales y correos \\
\texttt{ms-triage} & Episodios y priorización \\
\texttt{ms-analitica} & Tablero sobre MySQL sintético \\
\texttt{ms-ia-agente} & Agente conversacional \\
\bottomrule
\end{tabular}
\caption{Módulos incluidos en la ejecución.}
\end{table}

Las pruebas de contexto (\texttt{*ApplicationTests}, \texttt{EnvResolucionTest})
cargan la aplicación completa y necesitan el archivo \texttt{.env} y
PostgreSQL disponible. Las pruebas unitarias de seguridad y plantillas no
requieren base de datos.

\section{Herramientas}

\begin{itemize}
  \item JDK Eclipse Temurin 26.0.2 (el proyecto declara Java 17; el compilador
  usa el \textit{release} del proyecto).
  \item Maven Wrapper 3.9.16 incluido en cada módulo.
  \item Surefire para el resumen de pruebas.
  \item PowerShell 5+ para el script orquestador.
\end{itemize}

\section{Procedimiento}

\begin{enumerate}
  \item Verificar el stack de PostgreSQL y el archivo \texttt{.env}.
  \item Ejecutar, para cada módulo, \pantalla{.\mvnw.cmd -B test}
  desde su carpeta, con \texttt{JAVA\_HOME} apuntando al JDK disponible.
  \item Redirigir la salida completa a un log temporal y capturar el código de
  salida del proceso.
  \item Extraer la última línea de Surefire con el patrón
  \texttt{Tests run: ..., Failures: ..., Errors: ...}.
  \item Si el módulo falla, conservar las últimas 120 líneas del log en
  \pantalla{seguridad/resultados/maven/<modulo>-fallo.log}.
  \item Consolidar el resultado en \pantalla{resumen.csv} y
  \pantalla{resumen.md}.
\end{enumerate}

El script reproducible es
\pantalla{seguridad/pruebas/ejecutar-tests-maven.ps1}. Los logs completos se
guardan en la carpeta temporal del sistema para no inflar el repositorio.

\section{Criterios de aceptación}

Un módulo se marca \ok{} si el proceso termina con código 0 y Surefire no
reporta fallos ni errores. Cualquier código distinto de 0 se marca \fallo{} y
se conserva el final del log para el diagnóstico. No se usa
\texttt{-DskipTests} ni se excluye ninguna prueba.

\section{Limitaciones}

\begin{itemize}
  \item Se ejecutan las pruebas existentes; no se añadieron pruebas nuevas en
  esta fase (las funcionales y de seguridad están en documentos aparte).
  \item No hay reporte de cobertura: el proyecto no tiene JaCoCo configurado.
  \item No se usan Testcontainers: las pruebas de contexto dependen de la base
  local y del seed.
  \item La ejecución es secuencial por módulo; no se evaluó paralelismo.
  \item La primera ejecución puede tardar por descarga de dependencias.
\end{itemize}

\section{Reproducción}

\begin{verbatim}
powershell -NoProfile -ExecutionPolicy Bypass `
  -File seguridad/pruebas/ejecutar-tests-maven.ps1
\end{verbatim}

El script acepta \texttt{-Modules} para ejecutar un subconjunto, por ejemplo
\texttt{-Modules ms-seguridad,ms-triage}.

\end{document}
